% !TEX program = xelatex
\documentclass{英语试卷}
\begin{document}
\examtitle{2024 学年第二学期浙江省县域教研联盟高三年级模拟考试}{英语}
\par 英语试题\par
\par 命题：慈溪市教育局教研室 审题：长兴县教育研究中心 瑞安市教育发展研究院\par
\begin{notice}
\par 1. 答卷前，考生务必将自己的姓名、准考证号填写在答题卡上。\par
\par 2. 回答选择题时，选出每小题答案后，用铅笔把答题卡上对应题目的答案标号涂黑。如需改动，用橡皮擦干净后，再选涂其他答案标号。回答非选择题时，将答案写在答题卡上，写在本试卷上无效。\par
\par 3. 考试结束后，将本试卷和答题卡一并交回。微信公众号 IAI English\par
\end{notice}
\parttitle{第一部分 听力（共两节，满分30分）}
\par 做题时，先将答案标在试卷上，录音内容结束后，你将有两分钟的时间将试卷上的答案转涂到答题卡上。微信公众号IAI English\par
\sectitle{第一节（共5小题：每小题1.5分，满分7.5分）微信公众号IAI English}
\par 听下面5段对话。每段对话后有一个小题，从题中所给的A、B、C三个选项中选出最佳选项。听完每段对话后，你都有10秒钟的时间来回答有关小题和阅读下一小题。每段对话仅读一遍。\par
\question{1. What did the man do yesterday?}
\begin{choices}(3)
  \item He saw a film.
  \item He played football.
  \item He watched a game.
\end{choices}
\question{2. What are the speakers talking about?}
\begin{choices}(3)
  \item A secret plan.
  \item A royal castle.
  \item A chess match.
\end{choices}
\question{3. What's the relationship between the speakers?}
\begin{choices}(3)
  \item Professor and student.
  \item Employer and employee.
  \item Interviewer and interviewee.
\end{choices}
\question{4. What did the man just do?}
\begin{choices}(3)
  \item He called an ambulance. 微信公众号 IAI English
  \item He fell down from his bike.
  \item He jogged with the woman.
\end{choices}
\question{5. What are the speakers probably doing?}
\begin{choices}(3)
  \item Creating paintings.
  \item Exploring a museum.
  \item Taking an art history course.
\end{choices}
\sectitle{第二节（共15小题：每小题1.5分，满分22.5分}
\par 听下面5段对话或独白。每段对话或独白后有几个小题，从题中所给的A、B、C三个选项中选出最佳选项。听每段对话或独白前，你将有时间阅读各个小题，每小题5秒钟；听完后，各小题将给出5秒钟的作答时间。每段对话或独白读两遍。微信公众号IAI English\par
\par 听第6段材料，回答第6、7题。\par
\question{6. When will the speakers' party probably be held?}
\begin{choices}(3)
  \item In 18 days.
  \item In 20 days.
  \item In 22 days.
\end{choices}
\question{7. What are the speakers planning to do?}
\begin{choices}(3)
  \item Plant a tree.
  \item Clean the house.
  \item Decorate their home.
\end{choices}
\par 听第 7 段材料，回答第 8 至 9 题。\par
\question{8. Why does Geoffrey call Kate?}
\begin{choices}(3)
  \item To solve a question.
  \item To talk about their teacher.
  \item To help her with the homework.
\end{choices}
\question{9. What can we say about Kate?}
\begin{choices}(3)
  \item She's grateful.
  \item She's regretful.
  \item She's helpful.
\end{choices}
\par 听第8段材料，回答第10至12题。\par
\question{10. Where are the speakers probably?}
\begin{choices}(3)
  \item At their house.
  \item At a bus station.
  \item At a subway station.
\end{choices}
\question{11. What does the man suggest doing?}
\begin{choices}(3)
  \item Catching a bus.
  \item Calling a taxi.
  \item Waiting for a subway.
\end{choices}
\question{12. What is the woman concerned about?}
\begin{choices}(3)
  \item Her glasses.
  \item Her heavy bags.
  \item The distance to the airport.
\end{choices}
\par 听第9段材料，回答第13至16题。\par
\question{13. Why is the man being interviewed?}
\begin{choices}(3)
  \item He's publishing a book.
  \item He's acting in a movie.
  \item He's filming a TV series.
\end{choices}
\question{14. How did the man become successful?}
\begin{choices}(3)
  \item By studying at many schools.
  \item By doing many different jobs.
  \item By learning from a great coach.
\end{choices}
\question{15. What was the man's biggest challenge?}
\begin{choices}(3)
  \item Getting a stable job.
  \item Learning not to give up.
  \item Reading more books in the limited time.
\end{choices}
\question{16. What will the woman do next?}
\begin{choices}(3)
  \item Pass out some books.
  \item Talk to the next guest.
  \item Discuss her future plans.
\end{choices}
\par 听第10段材料，回答第17至20题。\par
\question{17. What is the talk mainly about?}
\begin{choices}(3)
  \item What to do in the park.
  \item How to get tickets for the park.
  \item How to win an amazing competition.
\end{choices}
\question{18. What can visitors do in the Space War Arena?}
\begin{choices}(3)
  \item Watch a show.
  \item Do a role-play.
  \item Get free coffee.
\end{choices}
\question{19. What event will happen tonight?}
\begin{choices}(3)
  \item A race.
  \item A concert.
  \item A theater show.
\end{choices}
\question{20. Who is the speaker probably?}
\begin{choices}(3)
  \item A hotel clerk.
  \item A café owner.
  \item A theme park worker.
\end{choices}
\parttitle{第二部分 阅读理解（共两节，满分50分）}
\sectitle{第一节（共15小题：每小题2.5分，满分37.5分）}
\par 阅读下列短文，从每题所给的A、B、C和D四个选项中，选出最佳答案。\par
\passagelabel{A}
\par While Iceland's aurora borealis(北极光) draws millions, this Nordic island offers year-round wonders far beyond it. From volcanic landscapes to sustainable tourism initiatives, here's why Iceland should top your travel list. The Golden Circle's Hidden Gems\par
\par Most visitors flock to the Golden Circle's iconic trio: Thingvellir National Park, Geysir hot springs, and Gullfoss waterfall. However, few explore the nearby Friðheimar Tomato Farm. This geothermal greenhouse serves tomato-based dishes(try the Bloody Mary soup!) while demonstrating Iceland's renewable energy innovations.\par
\par Westfjords: Untamed Wilderness 微信公众号 IAI English\par
\par Escape crowds in the remote Westfjords. Hike Hornstrandir Nature Reserve's cliffs where Arctic foxes roam freely. Don't miss the Dynjandi waterfall—a 100-meter waterfall resembling a bridal veil. Local tip: Rent a 4x4 car;\par
\par gravel roads require sturdy vehicles.\par
\par Sustainable Soaking\par
\par Blue Lagoon(泻湖) may be famous, but consider less-known geothermal pools. Sky Lagoon near Reykjavik offers a seven-step ritual blending Viking traditions with modern spa therapies. For eco-conscious travelers, Vok Baths in East Iceland use 100\% natural geothermal water without artificial heating.\par
\par Midnight Sun Adventures\par
\par Visit between May-July for 24-hour daylight. Join locals hiking Mount Esja at midnight or kayaking Jökulsárlón Glacier Lagoon under the golden glow. Pro photographers recommend capturing Seljalandsfoss waterfall with the sun hovering at 2 AM.\par
\par Responsible Travel Note\par
\par With overtourism concerns, authorities urge visitors to:\par
\par Stick to marked trails\par
\par Choose eco-certified accommodations (look for Vakinn labels)\par
\par Support rural communities by visiting East Iceland and the Highlands\par
\par When to Go\par
\par Winter(Nov-Mar): Northern Lights, ice caves\par
\par Summer(Jun-Aug): Midnight sun, puffin(海雀) colonies\par
\par Shoulder seasons: Lower prices, fewer crowds\par
\question{21. What is suggested for visitors exploring the Westfjords?}
\begin{choices}(2)
  \item Try tomato-based dishes.
  \item Watch out for Arctic foxes.
  \item Rent a 4x4 vehicle for gravel roads.
  \item Join the crowds at the Dynjandi waterfall.
\end{choices}
\question{22. Which activity is recommended for environmentally friendly travelers?}
\begin{choices}(2)
  \item Go on midnight sun adventures.
  \item Hike off marked trails in Hornstrandir.
  \item Soak in the Blue Lagoon's heated pools.
  \item Visit Vok Baths using natural geothermal water.
\end{choices}
\question{23. Where is the text probably taken from?}
\begin{choices}(2)
  \item A travel guide.
  \item A geography textbook.
  \item A product advertisement.
  \item An entertainment magazine.
\end{choices}
\passagelabel{B}
\par Dr. Aminata Diallo, a 52-year-old Malian engineer, never imagined her childhood chore — carrying water 8km daily — would inspire a revolution. Today, her solar-powered atmospheric water generators (AWGs) relieve the thirst of 300,000 Sahel villagers, earning her the 2023 Rolex Award for Enterprise. 微信公众号 IAI English\par
\par Born in drought-stricken Timbuktu, Diallo experienced her "eureka moment" during a sandstorm: "I watched camels lick condensation(冷凝水) off rocks and thought — why can't humans harvest air's moisture?" While studying mechanical engineering at MIT, she faced skepticism. "Professors said AWGs were impractical for deserts," she recalled. Determined, Diallo tested prototypes in her dorm shower's humidity, often working late into the night to refine her designs.\par
\par Traditional AWGs required expensive hygroscopic materials. Diallo's innovation? Treating bamboo fibers with sodium chloride — a technique inspired by Mali's salt caravans(商队). "Bamboo’s structure mimics camel nostrils (鼻孔) which efficiently take in moisture from the air," she explained at TEDGlobal 2022. Her \$50 portable units now produce 20 liters daily even at 15\% humidity, making clean water accessible to remote communities.\par
\par Diallo's nonprofit "Jiko Collective" trains women to assemble AWGs, changing traditional gender roles in conservative regions. "My first workshop had three grandmothers and a nursing mother," she told The Guardian, "they fit machine parts while watching babies, proving that innovation thrives when everyone contributes." Since2020, over 4,000 women have learned these skills, giving clean water to 300,000 people and transforming their communities' social fabric.\par
\par When tech giants offered millions for her patent, Diallo refused. "Water belongs to the sky, not shareholders," she declared. Instead, she open-sourced designs through UNESCO, making them available for worldwide adoption. Critics call it naive, but to this award-winning engineer, every water drop from air shouts, "The sky's gift belongs to all."\par
\question{24. What initially motivated Diallo's invention?}
\begin{choices}(2)
  \item Academic pressure at MIT.
  \item Childhood experiences in Mali.
  \item Encouragement from professors.
  \item Inspiration from camel behavior.
\end{choices}
\question{25. What does the underlined word "hygroscopic" in paragraph 3 most likely mean?}
\begin{choices}(2)
  \item Water-absorbing.
  \item Energy-saving.
  \item Heat-resistant.
  \item Cost-effective.
\end{choices}
\question{26. Why did Diallo refuse tech companies' offers?}
\begin{choices}(2)
  \item To focus on academic research.
  \item To protect her technical secrets.
  \item To reduce production costs further.
  \item To keep designs publicly accessible.
\end{choices}
\question{27. Which of the following best describe Diallo's character?}
\begin{choices}(2)
  \item Ambitious and profit-driven.
  \item Competitive and self-centered.
  \item Observant and socially conscious.
  \item Conservative and strong-minded.
\end{choices}
\passagelabel{C}
\par Once dismissed as a mental interference, daydreaming is undergoing a scientific reevaluation. In our productivity-obsessed era, where multitasking reduces attention spans to me re seconds, this once-criticized mental state is now recognized as a vital cognitive process. Groundbreaking studies from leading universities reveal its multifaceted benefits, challenging conventional wisdom about focus and creativity. 微信公众号 IAI English\par
\par Harvard neuroscientists recently demonstrated the mechanics behind this phenomenon. Their 2024 study in Cognitive Psychology had 150 participants complete creative problem-solving tasks. Those practicing structured daydreaming — guiding thoughts around specific themes like "future inventions"—outperformed intensely focused peers by 30\%. Brain imaging revealed heightened activity in the default mode network (DMN), a neural circuit previously mislabeled as idle during rest. "The DMN acts as a mental workshop," explains lead researcher Dr. Laura Bennett. "When consciously directed, it connects distant concepts — like linking childhood memories to workplace challenges — to spark innovation."\par
\par Another study from Stanford University highlights daydreaming's role in emotional regulation. In controlled trials, subjects who daydreamed for 10 minutes after high-stress simulations showed 25\% lower cortisol levels and improved emotional clarity. "This isn't escapism," clarifies psychologist Dr. Emily Carter. "It's the brain's self-maintenance mode, converting chaotic emotions into manageable narratives."\par
\par Interestingly, daydreaming also improves long-term memory. A 2023 experiment at the University of Cambridge showed that students who daydreamed during breaks between study sessions retained 20\% more information than those who didn't. "This suggests that daydreaming helps consolidate learning," notes Dr. James Wilson.\par
\par Despite these benefits, experts caution against excessive daydreaming. It is indicated that excessive, undirected daydreaming correlates with a 17\% increase in reality-disconnection symptoms among adolescents. "While beneficial in modernation, too much daydreaming can lead to disengagement from reality," warns Dr.Bennett. Balancing focused work with periods of mindful daydreaming is key to harnessing its full potential.\par
\question{28. What was the conventional view of daydreaming?}
\begin{choices}(2)
  \item A vital cognitive process.
  \item A distraction of mental focus.
  \item A way to improve memory.
  \item A method to enhance creativity.
\end{choices}
\question{29. What can be inferred about the brain's default mode network(DMN)?}
\begin{choices}(2)
  \item It remains inactive during both rest and focused tasks.
  \item It only activates during stressful situations to reduce anxiety.
  \item It was once considered inactive but now known to link ideas.
  \item It primarily strengthens long-term memory through conscious effort.
\end{choices}
\question{30. Which outcome is associated with structured daydreaming in the text?}
\begin{choices}(2)
  \item Disengagement from reality.
  \item Decreased long-term memory.
  \item Strengthened stress and anxiety.
  \item Enhanced emotional management.
\end{choices}
\question{31. Which real-life situation best matches the research findings on daydreaming?}
\begin{choices}(2)
  \item A student avoids all breaks to maintain focus.
  \item An artist schedules deliberate "mind-wandering time".
  \item A teacher bans daydreaming in class to ensure discipline.
  \item A programmer works 12-hour shifts without mental pauses.
\end{choices}
\passagelabel{D}
\par In a quiet corner of Oxford University's Bodleian Library, a robotic arm flips through a 15th-century manuscript, its high-resolution cameras capturing every faded in k stroke. This is ARCHiVE (Automated Robotic Cultural Heritage Inspection and Verification Engine), a £2.3 million project revolutionizing how fragile artifacts are preserved and studied. Combining robotics, artificial intelligence, and advanced imaging, ARCHiVE represents a model shift in cultural conservation.\par
\par Traditional conservation methods often risk damaging delicate materials. ARCHiVE's robotic system, equipped with micro-sensors and AI algorithms, can detect cracks invisible to the human eye and predict structural weaknesses before they worsen. "It's like having a doctor for ancient books," explains Dr. Emily Carter, lead conservator. Since its 2022 launch, ARCHiVE has scanned over 8,000 items, from Egyptian papyri to medieval maps, with 99.7\% accuracy in damage assessment. Its non-invasive approach has reduced physical handling by 85\% , significantly lowering the risk of accidental damage.\par
\par The project's digital twin technology creates exact virtual copies, allowing global access to restricted collections. During the 2023 Cairo Museum flood, ARCHiVE's digital models enabled researchers to virtually "restore" water-damaged artifacts within hours. "We’re not just saving objects," says Carter, "we’re preserving humanity’s collective memory." This technology has also democratized access, with over 50,000 digital artifacts now available to researchers and the public worldwide. 微信公众号 IAI English\par
\par However, some scholars argue that over-reliance on technology may distance researchers from hands-on experience. "You can't smell history through a screen," protests Prof. James Whitaker, a traditionalist at Cambridge. Others worry about data ownership, as scanned artifacts become digital assets potentially controlled by tech corporations. A 2024 survey revealed that 62\% of conservators fear losing control over digitized cultural heritage.\par
\par ARCHiVE's next phase involves 4D scanning, capturing how materials change over time. Pilot tests on Viking ship timbers have already revealed previously unknown construction techniques. Meanwhile, UNESCO is developing global standards for digital heritage, ensuring these technological marvels serve humanity rather than corporate interests. As Carter reflects, "Our ancestors left us these treasures; our duty is to pass them forward—both physically and digitally."\par
\question{32. What is the main function of ARCHiVE's robotic system?}
\begin{choices}(2)
  \item Repairing damaged artifacts.
  \item Creating virtual museum tours.
  \item Identifying and forecasting damage.
  \item Selling digital copies of artifacts.
\end{choices}
\question{33. How did ARCHiVE help during the Cairo Museum flood?}
\begin{choices}(2)
  \item By raising emergency funds online.
  \item By training local conservation staff.
  \item By physically moving artifacts to safety.
  \item By providing digital models for restoration.
\end{choices}
\question{34. What is Prof. Whitaker's main concern about ARCHiVE?}
\begin{choices}(2)
  \item Data ownership concerns.
  \item Loss of direct experience.
  \item High operating costs.
  \item Environmental impact.
\end{choices}
\question{35. Which of the following would be the best title for the article?}
\begin{choices}(2)
  \item ARCHiVE: Digital Guardians of the Past.
  \item Heritage Conservation: Traditional or Digital.
  \item Human Conservators: Why They Still Matter.
  \item AI in Museums: Ethical Dilemmas and Solutions.
\end{choices}
\sectitle{第二节（共5小题；每小题2.5分，满分12.5分）}
\par 阅读下面短文，从短文后的选项中选出能填入空白处的最佳选项。选项中有两项为多余选项。\par
\par The Enduring Power of Literature\par
\par In an age dominated by digital media, the question of literature's relevance has never been more pressing. \blanknum{36} As a novelist and literary critic, I often turn to the insights of great writers who have delved into this question before me.\par
\par Toni Morrison once said, "We die. That may be the meaning of life. But we do language. That may be the measure of our lives." For Morrison, literature was not merely a form of entertainment. \blanknum{37} Through language, we preserve our experiences, our emotions, and our very essence. 微信公众号 IAI English\par
\par James Baldwin, in his essay "The Creative Process", argued that literature serves as a mirror and a map. He wrote, “\blanknum{38}” Baldwin believed that literature had the power to confront uncomfortable truths and inspire social change. His works demonstrate how stories can simultaneously reflect reality and chart moral progress. The novel Go Tell It on the Mountain, for instance, not only exposed racial injustice in 1950s America but also envisioned paths toward empathy and reconciliation(和解).\par
\par Milan Kundera, in "The Art of the Novel", emphasized literature's ability to explore the complexity of human existence. "The novel's wisdom," he claimed, "is different from the wisdom of philosophy. \blanknum{39}" For Kundera, literature's value lies in its capacity to embrace life's contradictions without seeking simplistic resolutions.\par
\par As I reflect on these perspectives, I am reminded of literature's unique role in our lives. \blanknum{40} By revealing unseen possibilities, it acts as both a time capsule of collective memory and a blueprint for unwritten futures. In a world increasingly driven by binary(二进制) thinking, literature offers a space for subtle difference, empathy, and profound understanding.\par
\begin{choices}(1)
  \item It is an act of resistance against being forgotten.
  \item Literature's ability to entertain is its most important function.
  \item Literature provides clear answers to life's most difficult questions.
  \item Some even argue that literature has lost its place in modern society.
  \item It is not just a reflection of who we are but also a vision of who we might become.
  \item It is the wisdom of uncertainty, of ambiguity and of questioning rather than answering.
  \item It was through literature that we could see ourselves clearly and chart a course for our future.
\end{choices}
\parttitle{第三部分 语言知识运用（共两节，满分30分）}
\sectitle{第一节（共15小题；每小题1分，满分15分）}
\par 阅读下面短文，从短文后各题所给的A、B、C和D四个选项中，选出可以填入空白处的最佳选项。\par
\par During a night shift at the vet clinic, my colleague mentioned a Yorkshire Terrier puppy. It had treatable diarrhoea but was \blanknum{41} for euthanasia(安乐死) simply because its owner refused to pay for the \blanknum{42}. Heartbroken, I called the owner, convincing him to \blanknum{43} the puppy to me. I treated him, named him Atticus, and he has thrived as my close companion for eight years. Sadly, not all dogs are as \blanknum{44}.\par
\par After becoming a vet in 2015, I witnessed countless pets brought in for euthanasia because their owners no longer wanted them. \blanknum{45} ranged from financial struggles to housing restrictions or behavioral issues. Determined to help, I began taking \blanknum{46} animals home. Word spread, and before long, I'd \blanknum{47} a variety of dogs and cats. Eventually, with the support of friends who also became foster \blanknum{48}, I launched a non-profit organization—Celandine Wood Animal Rescue. The cost-of-living crisis saw a(an) \blanknum{49} in animals being given away.\par
\par Then came Dave, an eight-month-old Cane Corso. He’d been imported for breeding, but was found out with elbow dysplasia(肘关节发育不良). Considering him \blanknum{50}, the owner gave him up. Our \blanknum{51} saved him and found him a wonderful home. Currently, we care for about 50 dogs, 40 cats, and wildlife like birds and foxes.\par
\par Stories like this \blanknum{52} our mission, yet countless animals still slip through the cracks. \blanknum{53}, we're fundraising to build a larger shelter, providing better resources and space. We never put an animal down unless we've \blanknum{54} all other options. 微信公众号 IAI English\par
\par Every life matters. 55 there are animals like Atticus and Dave, we'll keep fighting for them.\par
\cz{41}{evaluated}{scheduled}{qualified}{subscribed}
\cz{42}{examination}{expense}{surgery}{treatment}
\cz{43}{hand down}{sign over}{set aside}{put away}
\cz{44}{fortunate}{healthy}{curable}{adorable}
\cz{45}{Reasons}{Conflicts}{Results}{Occasions}
\cz{46}{abused}{innocent}{dying}{abandoned}
\cz{47}{accumulated}{adapted}{discovered}{drawn}
\cz{48}{workers}{partners}{carers}{trainers}
\cz{49}{effort}{change}{trend}{event}
\cz{50}{aggressive}{disabled}{useless}{unprofitable}
\cz{51}{camp}{club}{shelter}{project}
\cz{52}{extend}{fuel}{guide}{reflect}
\cz{53}{Obviously}{Consequently}{Particularly}{Immediately}
\cz{54}{exhausted}{managed}{investigated}{avoided}
\cz{55}{Now that}{Even if}{As long as}{For fear that}
\sectitle{第二节（共10小题；每小题1.5分，满分15分）}
\par 阅读下面短文，在空白处填入1个适当的单词或括号内单词的正确形式。\par
\par An unexpected cultural phenomenon has sparked widespread discussions among netizens recently. Some users shared posts\blanknum{56}(claim) that while visiting relatives in the UK, they discovered British children using Chinese "Hengshui-style" English writing practice books. 微信公众号 IAI English\par
\par Hengshui Style, also known as "Hengshui Ti", is a standardized English handwriting method popularized by students from Hengshui High School \blanknum{} strictness and academic excellence are well known in China. The writing style stands out for its \blanknum{} (discipline) structure, neat formatting, and aesthetically pleasing appearance, resembling printed text. Over the years, it has earned a reputation as \blanknum{} "secret weapon for exams", which is believed by many educators and students \blanknum{} (be) highly beneficial for improving legibility and efficiency during timed tests.\par
\par The online buzz triggered nostalgia among Chinese netizens. Many \blanknum{61}(fond) recalled their own school days, carefully practicing English sentences in "Hengshui-style" workbooks, some even still possessing identical copybooks. Social media platforms\blanknum{62}(flood) with users sharing their own "Hengshui-style" handwriting samples and offering tutorials and tips for mastering the style. Some shared before-and-after \blanknum{63}(compare), showcasing dramatic improvements in their penmanship, \blanknum{64} others humorously debated whether the strict style limited creativity.\par
\par This cross-cultural fascination highlights the global reach of educational practices and the enduring legacy of a handwriting method once \blanknum{65} (confine) to Chinese classrooms.\par
\parttitle{第四部分 写作（共两节，满分40分）}
\sectitle{第一节（满分15分）}
\par 假设你是李华，你校学生会正在为“高三励志墙”征集英文标语。请你给负责此次活动的外教 Mr. Smith 写一封邮件，内容包括：(1) 你的标语；(2) 推荐理由。\par
\par 注意：微信公众号IAI English\par
\par (1) 写作词数应为 80 左右;\par
\par (2) 可适当增加细节，以使行文连贯。\par
\sectitle{第二节（满分25分）}
\par 阅读下面材料，根据其内容和所给段落开头语续写两段，使之构成一篇完整的短文。\par
\par Amanda and Susie, though sisters with a two-year age gap, had never got along. From the minute Susie was born, Amanda knew the baby disliked her.\par
\par As Susie grew, she hated Amanda for telling her how to do things. Amanda viewed Susie as spoilt, because she was the youngest in the family. At secondary school, Amanda shone at English; Susie was a maths whizz. Each made the other feel inadequate. When Amanda turned fifteen, Susie was angry when her sister was allowed to go to the youth club disco. Perhaps that was why Susie married so young — anything to beat her sister down the aisle. But Amanda got her revenge. When it was her turn, she chose a bridesmaids clashed with Susie's hair at her wedding.\par
\par Now, despite being in their fifties, these grievances(牢骚) were still fresh in the sisters' minds. The only thing they had in common was that they both loved their mother deeply.\par
\par Then one day, her father rang with the terrible news that their mother had died suddenly. Both sisters were stricken with grief, but Susie, was also upset that their father hadn't informed her first.\par
\par Then their father told them about their mother's will. "She’s left you the painting." He said to them both. The painting — a portrait of the sisters by the mother — was the source of one of their biggest childhood arguments. To their surprise, it was bequeathed(遗赠) to both, "Your mom hopes that you each have it for six months and then swap over, for all your rest life." added the father. 微信公众号 IAI English\par
\par "We live miles apart. We could get it delivered." said Susie. "It's against the will," their father replied, "'I love to look at it, especially now you've both left home.' Your mom used to say. Now it felt like the last link to her."\par
\par For the sake of their beloved mother, the two sister began the cyclic keeping of the painting. This time it was Susie's turn to deliver the painting. 注意：微信公众号 IAI English\par
\par (1) 续写词数应为 150 左右;\par
\par (2) 请按如下格式在答题卡的相应位置作答。\par
\writingpara{Paragraph 1: When Susie arrived, she found Amanda lying in the bed with her leg broken in an accident.}
\writinglines{8}
\writingpara{Paragraph 2: Silently, they began to admire the painting, and all of a sudden, a line of little words beneath draw Amanda's attention.}
\writinglines{8}
\end{document}
